% ============================================================
% AgroSmart Andino -- Guia de Implementacion Completa
% Pasos 1 al 8: Local + AWS
% ============================================================
\documentclass[11pt,a4paper]{article}

\usepackage[utf8]{inputenc}
\usepackage[T1]{fontenc}
\usepackage[spanish,es-nodecimaldot,es-noindentfirst]{babel}
\usepackage[margin=2cm,top=2.5cm,bottom=2.5cm]{geometry}
\usepackage{hyperref}
\usepackage{xcolor}
\usepackage{booktabs}
\usepackage{array}
\usepackage{listings}
\usepackage{tcolorbox}
\usepackage{enumitem}
\usepackage{fancyhdr}
\usepackage{titlesec}
\usepackage{multicol}
\usepackage{amsmath}
\usepackage{amssymb}

% Paleta
\definecolor{agroverde}{RGB}{27,67,50}
\definecolor{agrolight}{RGB}{45,106,79}
\definecolor{agroclear}{RGB}{184,233,199}
\definecolor{agroazul}{RGB}{0,80,180}
\definecolor{agrorojo}{RGB}{180,30,30}
\definecolor{agronaranja}{RGB}{200,100,0}
\definecolor{agrogris}{RGB}{90,90,90}
\definecolor{codebg}{RGB}{245,245,245}

% Titulos
\titleformat{\section}
  {\large\bfseries\color{agroverde}}
  {\thesection.}{0.5em}{}
  [\color{agroverde}\hrule]
\titleformat{\subsection}
  {\normalsize\bfseries\color{agroazul}}
  {\thesubsection}{0.5em}{}

% Header/Footer
\pagestyle{fancy}
\fancyhf{}
\fancyhead[L]{\small\color{agrogris} AgroSmart Andino -- Gu\'ia de Implementaci\'on}
\fancyhead[R]{\small\color{agrogris} Maestr\'ia IA -- UNI -- Big Data 2026}
\fancyfoot[C]{\small\thepage}
\renewcommand{\headrulewidth}{0.4pt}

% Listings
\lstset{
  basicstyle=\ttfamily\small,
  backgroundcolor=\color{codebg},
  frame=single, framerule=0pt,
  breaklines=true, showstringspaces=false,
  keywordstyle=\color{agroazul}\bfseries,
  commentstyle=\color{agrogris}\itshape,
  stringstyle=\color{agrorojo},
  numbers=none,
  literate=
    {\'a}{{\'{a}}}1 {\'e}{{\'{e}}}1 {\'i}{{\'{i}}}1
    {\'o}{{\'{o}}}1 {\'u}{{\'{u}}}1 {\~n}{{\~{n}}}1
    {\'A}{{\'{A}}}1 {\'E}{{\'{E}}}1 {\'I}{{\'{I}}}1
    {\'O}{{\'{O}}}1 {\'U}{{\'{U}}}1
}
\lstdefinestyle{bash}{language=bash,
  morekeywords={cd,python,streamlit,pip,conda,cqlsh,docker},
  basicstyle=\ttfamily\small,
  backgroundcolor=\color{black!90},
  basicstyle=\ttfamily\small\color{white},
  keywordstyle=\color{yellow},
  commentstyle=\color{gray},
  frame=none,
}
\lstdefinestyle{python}{language=python,
  basicstyle=\ttfamily\scriptsize,
  backgroundcolor=\color{codebg},
  keywordstyle=\color{agroazul}\bfseries,
  commentstyle=\color{agrogris}\itshape,
  stringstyle=\color{agrorojo},
}

% Cajas de colores
\tcbuselibrary{skins}

\newtcolorbox{localbox}[1][Ejecuci\'on LOCAL]{
  colback=agroazul!8, colframe=agroazul,
  title=#1, fonttitle=\bfseries\color{white},
  colbacktitle=agroazul, arc=4pt, boxrule=0.5pt,
}
\newtcolorbox{awsbox}[1][Ejecuci\'on AWS]{
  colback=agronaranja!8, colframe=agronaranja!80!black,
  title=#1, fonttitle=\bfseries\color{white},
  colbacktitle=agronaranja!80!black, arc=4pt, boxrule=0.5pt,
}
\newtcolorbox{alertbox}[1][Importante]{
  colback=agrorojo!8, colframe=agrorojo,
  title=#1, fonttitle=\bfseries\color{white},
  colbacktitle=agrorojo, arc=4pt, boxrule=0.5pt,
}
\newtcolorbox{okbox}[1][Resultado esperado]{
  colback=agroclear!50, colframe=agrolight,
  title=#1, fonttitle=\bfseries\color{white},
  colbacktitle=agrolight, arc=4pt, boxrule=0.5pt,
}

% ============================================================
\begin{document}
% ============================================================

\begin{titlepage}
  \centering
  \vspace*{1cm}
  {\color{agroverde}\rule{\linewidth}{2pt}}\\[0.5em]
  {\huge\bfseries\color{agroverde} AgroSmart Andino}\\[0.3em]
  {\large\color{agroverde} Gu\'ia de Implementaci\'on Completa}\\[0.3em]
  {\normalsize\color{agrogris} Pasos 1 al 8 --- Local y AWS paso a paso}\\[0.3em]
  {\color{agroverde}\rule{\linewidth}{2pt}}

  \vspace{1.5cm}
  \begin{tcolorbox}[colback=agroverde!5, colframe=agroverde,
    arc=6pt, boxrule=1pt]
    \begin{center}
      \textbf{Pipeline Big Data con Dashboard en Tiempo Real}\\[0.3em]
      \normalsize
      Monitoreo de Papa Nativa en Huangascar, Yauyos, Lima\\[0.5em]
      \begin{tabular}{ll}
        \textbf{Curso:} & Big Data\\
        \textbf{Docente:} & Mg. Rosa Virginia Encinas Quille\\
        \textbf{Universidad:} & Universidad Nacional de Ingenier\'ia (UNI)\\
        \textbf{Programa:} & Maestr\'ia en Inteligencia Artificial\\
        \textbf{Grupo:} & 2\\
        \textbf{Integrantes:} & Montes E. $\cdot$ Vizcardo F. $\cdot$ Massa Q.
                                $\cdot$ Huali C. $\cdot$ Lazaro G.\\
        \textbf{Fecha:} & Mayo 2026\\
      \end{tabular}
    \end{center}
  \end{tcolorbox}

  \vspace{1.5cm}
  \textbf{Resumen del sistema:}\\[0.4em]
  \begin{tabular}{llll}
    \toprule
    Componente & Tecnolog\'ia & Registros & Estado \\
    \midrule
    Data Acquisition & Python 3.13 & 72\,397 & {\color{agrolight}\checkmark} \\
    Data Lake & AWS S3 & 3 CSV & {\color{agrolight}\checkmark} \\
    ETL + ML & Spark 3.4 / EMR & 3\,653 Gold & {\color{agrolight}\checkmark} \\
    Indexaci\'on & 4 \'indices & 18/47/31 d\'ias & {\color{agrolight}\checkmark} \\
    NoSQL & Cassandra 4.1 & 4\,459 dry-run & {\color{agrolight}\checkmark} \\
    Dashboard & Streamlit 1.x & 11 KPIs & {\color{agrolight}\checkmark} \\
    Informe PDF & LaTeX & 20 p\'aginas & {\color{agrolight}\checkmark} \\
    Presentaci\'on & Beamer & 17 diapositivas & {\color{agrolight}\checkmark} \\
    \bottomrule
  \end{tabular}

  \vfill
  {\color{agrogris}\small Directorio base del proyecto:
    \texttt{d:\textbackslash{}Documents\textbackslash{}GitHub\textbackslash{}Maestria\_IA\_caso\_practico\textbackslash{}Big\_data\textbackslash{}trabajo\_final\textbackslash{}agrosmart\_andino\textbackslash{}}}
\end{titlepage}

\tableofcontents
\newpage

% ============================================================
\section{Prerrequisitos e Instalaci\'on del Entorno}
% ============================================================

\subsection{Software necesario (instalaci\'on local Windows)}

\begin{tabular}{lll}
\toprule
Software & Versi\'on probada & Instalaci\'on \\
\midrule
Python (Miniconda) & 3.13 & \href{https://docs.conda.io}{conda install} \\
MiKTeX & 25.12 & \href{https://miktex.org}{miktex.org} (para PDF) \\
AWS CLI & 2.x & \href{https://aws.amazon.com/cli/}{aws.amazon.com/cli} \\
Docker Desktop & 24.x & Para Cassandra local (opcional) \\
VS Code & cualquiera & Editor recomendado \\
\bottomrule
\end{tabular}

\subsection{Instalar dependencias Python}

\begin{localbox}[Instalar paquetes Python]
\begin{lstlisting}[style=bash]
# Desde la carpeta agrosmart_andino/
conda activate base     # o el entorno de tu preferencia
pip install -r requirements.txt
\end{lstlisting}
\end{localbox}

Contenido de \texttt{requirements.txt}:
\begin{lstlisting}[style=python]
pandas>=2.0
numpy>=1.24
requests>=2.31
boto3>=1.34
python-dotenv>=1.0
streamlit>=1.30
plotly>=5.18
cassandra-driver>=3.29
pyspark>=3.4      # solo necesario si ejecutas localmente sin EMR
\end{lstlisting}

\subsection{Estructura de directorios}

\begin{lstlisting}
agrosmart_andino/
+-- data/
|   +-- raw/           <- CSVs descargados (Parte 2)
|   +-- gold/          <- CSVs procesados para dashboard (Parte 4+5)
+-- scripts/
|   +-- data_acquisition/
|   |   +-- 01_download_nasa_power.py
|   |   +-- 02_generate_iot_data.py
|   |   +-- 03_prepare_midagri_data.py
|   |   +-- 04_validate_datasets.py
|   +-- aws/
|   |   +-- 05_upload_to_s3.py
|   |   +-- 06_create_emr_cluster.py
|   +-- cassandra/
|       +-- 07_create_schema.cql
|       +-- 08_load_data.py
|       +-- 09_setup_cassandra_cloud9.sh
+-- notebooks/
|   +-- agrosmart_pipeline.ipynb   <- Spark ETL + ML (38 celdas)
+-- dashboard/
|   +-- app.py                     <- Streamlit dashboard
+-- docs/
|   +-- informe_tecnico.tex/.pdf   <- Informe (20 pags)
|   +-- presentacion.tex/.pdf      <- Diapositivas (17 slides)
|   +-- compilar.bat
+-- .env                           <- Credenciales AWS (NO subir a git)
+-- .env.example
+-- requirements.txt
+-- README.md
\end{lstlisting}

% ============================================================
\section{Parte 1 -- Estructura del Proyecto}
% ============================================================

\textbf{Objetivo:} Crear la estructura de carpetas y archivos de configuraci\'on base.

\begin{localbox}[Crear la estructura de carpetas]
\begin{lstlisting}[style=bash]
# Crear directorios desde trabajo_final/
mkdir -p agrosmart_andino/data/raw
mkdir -p agrosmart_andino/data/gold
mkdir -p agrosmart_andino/scripts/data_acquisition
mkdir -p agrosmart_andino/scripts/aws
mkdir -p agrosmart_andino/scripts/cassandra
mkdir -p agrosmart_andino/notebooks
mkdir -p agrosmart_andino/dashboard
mkdir -p agrosmart_andino/docs
\end{lstlisting}
\end{localbox}

\begin{localbox}[Configurar archivo .env]
\begin{lstlisting}[style=bash]
# Copiar plantilla y editar con tus credenciales AWS
copy agrosmart_andino\.env.example agrosmart_andino\.env
# Editar .env y completar:
#   AWS_ACCESS_KEY_ID=ASIA...
#   AWS_SECRET_ACCESS_KEY=...
#   AWS_SESSION_TOKEN=...
#   S3_BUCKET_NAME=agrosmart-andino-demo
#   AWS_DEFAULT_REGION=us-east-1
\end{lstlisting}
\end{localbox}

\begin{okbox}[Resultado Parte 1]
Estructura de carpetas creada. Archivos: README.md, requirements.txt, .gitignore, .env.example.
\end{okbox}

% ============================================================
\section{Parte 2 -- Adquisici\'on y Validaci\'on de Datos}
% ============================================================

\textbf{Objetivo:} Descargar y generar los 3 datasets fuente. Validar 72\,397 registros.

\subsection{Script 01 -- Descarga NASA POWER API}

\begin{localbox}[Ejecutar desde agrosmart\_andino/]
\begin{lstlisting}[style=bash]
cd agrosmart_andino
python scripts/data_acquisition/01_download_nasa_power.py
\end{lstlisting}
\end{localbox}

\textbf{Qu\'e hace:} Llama la API gratuita de NASA POWER
(\texttt{https://power.larc.nasa.gov/api/temporal/daily/}) para 5 estaciones
en Yauyos. Descarga 10 a\~nos de datos diarios (2015--2024).

\begin{okbox}[Resultado esperado]
\texttt{data/raw/nasa\_power\_yauyos.csv} --- 18\,265 filas, 1.6 MB\\
5 estaciones: Huangascar, Yauyos, Carania, Tanta, Vilca
\end{okbox}

\begin{alertbox}[Si hay error de red]
La API NASA POWER es gratuita y no requiere clave. Si hay timeout, ejecutar
nuevamente. La URL base es:
\texttt{https://power.larc.nasa.gov/api/temporal/daily/point}
\end{alertbox}

\subsection{Script 02 -- Generar datos IoT sint\'eticos}

\begin{localbox}[]
\begin{lstlisting}[style=bash]
python scripts/data_acquisition/02_generate_iot_data.py
\end{lstlisting}
\end{localbox}

\textbf{Qu\'e hace:} Genera datos horarios de 3 nodos IoT simulados para
las coordenadas de Huangascar. Usa distribuciones calibradas con el clima
real de la zona.

\begin{okbox}[]
\texttt{data/raw/sensores\_iot.csv} --- 52\,632 filas, 4.6 MB\\
3 nodos: NODO-001 (3\,000m), NODO-002 (3\,200m), NODO-003 (3\,500m)
\end{okbox}

\subsection{Script 03 -- Preparar datos MIDAGRI}

\begin{localbox}[]
\begin{lstlisting}[style=bash]
python scripts/data_acquisition/03_prepare_midagri_data.py
\end{lstlisting}
\end{localbox}

\begin{okbox}[]
\texttt{data/raw/produccion\_papa\_yauyos.csv} --- 1\,500 filas\\
Producci\'on hist\'orica de papa nativa por distrito y variedad (2000--2024)
\end{okbox}

\subsection{Script 04 -- Validar los 3 datasets}

\begin{localbox}[]
\begin{lstlisting}[style=bash]
python scripts/data_acquisition/04_validate_datasets.py
\end{lstlisting}
\end{localbox}

\textbf{Qu\'e valida:} Existencia, filas m\'inimas, duplicados, nulos cr\'iticos,
rangos de variables (temperatura, humedad, precipitaci\'on).

\begin{okbox}[]
\begin{lstlisting}
[OK] nasa_power_yauyos.csv    -- 18,265 filas  -- VALIDO
[OK] sensores_iot.csv         -- 52,632 filas  -- VALIDO
[OK] produccion_papa_yauyos.csv -- 1,500 filas -- VALIDO
Total: 72,397 registros validados
\end{lstlisting}
\end{okbox}

% ============================================================
\section{Parte 3 -- Data Lake en AWS S3}
% ============================================================

\textbf{Objetivo:} Subir los 3 CSVs validados a un bucket S3 como Data Lake.

\subsection{Paso previo: Crear el bucket S3 en AWS Console}

\begin{awsbox}[AWS Console -- Crear bucket S3]
\begin{enumerate}
  \item Iniciar sesi\'on en \href{https://awsacademy.instructure.com}{AWS Academy}
  \item Abrir el Learner Lab $\to$ \textbf{[Start Lab]} $\to$ esperar indicador VERDE
  \item Clic en \textbf{[AWS]} para abrir la consola
  \item Buscar \textbf{S3} $\to$ \textbf{Create bucket}
  \item Configurar:
  \begin{itemize}
    \item Bucket name: \texttt{agrosmart-andino-demo}
    \item Region: \textbf{US East (N. Virginia) us-east-1}
    \item Block Public Access: activado (dejar por defecto)
    \item Versioning: Disabled
  \end{itemize}
  \item Clic \textbf{Create bucket}
\end{enumerate}
\end{awsbox}

\subsection{Copiar credenciales al .env}

\begin{awsbox}[AWS Academy -- Copiar credenciales]
\begin{enumerate}
  \item En Learner Lab $\to$ \textbf{AWS Details}
  \item Copiar los 3 valores y pegarlos en \texttt{.env}:
\end{enumerate}
\begin{lstlisting}
AWS_ACCESS_KEY_ID=ASIA...
AWS_SECRET_ACCESS_KEY=...
AWS_SESSION_TOKEN=IQoJ...  <- token largo, NO OMITIR
S3_BUCKET_NAME=agrosmart-andino-demo
AWS_DEFAULT_REGION=us-east-1
\end{lstlisting}
\end{awsbox}

\begin{alertbox}[Las credenciales de AWS Academy expiran en 3 horas]
Cada vez que reinicies el lab debes copiar nuevas credenciales. El
\texttt{AWS\_SESSION\_TOKEN} es obligatorio; sin \'el la conexi\'on falla.
\end{alertbox}

\subsection{Ejecutar la subida}

\begin{localbox}[Desde agrosmart\_andino/]
\begin{lstlisting}[style=bash]
python scripts/aws/05_upload_to_s3.py
\end{lstlisting}
\end{localbox}

\begin{okbox}[]
\begin{lstlisting}
[OK] s3://agrosmart-andino-demo/raw/clima/nasa_power_yauyos.csv
[OK] s3://agrosmart-andino-demo/raw/sensores/sensores_iot.csv
[OK] s3://agrosmart-andino-demo/raw/produccion/produccion_papa_yauyos.csv
Total almacenado: 6.3 MB
\end{lstlisting}
\end{okbox}

% ============================================================
\section{Parte 4 -- Pipeline ETL y ML con Apache Spark en EMR}
% ============================================================

\textbf{Objetivo:} Crear cl\'uster EMR, ejecutar el notebook PySpark con
ETL + \'indices agroclim\'aticos + modelos ML.

\subsection{Crear el cl\'uster EMR}

\begin{awsbox}[Opci\'on A: Script autom\'atico]
\begin{lstlisting}[style=bash]
# Desde agrosmart_andino/ (necesita credenciales en .env)
python scripts/aws/06_create_emr_cluster.py
# Esperar ~8-10 minutos hasta estado WAITING
\end{lstlisting}
\end{awsbox}

\begin{awsbox}[Opci\'on B: AWS Console manualmente]
\begin{enumerate}
  \item AWS Console $\to$ \textbf{EMR} $\to$ \textbf{Create cluster}
  \item Configuraci\'on:
  \begin{itemize}
    \item Nombre: \texttt{AgroSmart-Cluster}
    \item EMR release: \textbf{emr-6.15.0}
    \item Applications: \textbf{Spark, JupyterHub}
    \item Primary (master): \texttt{m5.xlarge}
    \item Core: \texttt{m5.xlarge} (1 nodo)
    \item Key pair: crear o usar existente
  \end{itemize}
  \item Clic \textbf{Create cluster} $\to$ esperar estado \textbf{Waiting}
\end{enumerate}
\end{awsbox}

\subsection{Conectarse a JupyterHub y ejecutar el notebook}

\begin{awsbox}[Acceder a JupyterHub en EMR]
\begin{enumerate}
  \item EMR Console $\to$ tu cl\'uster $\to$ pesta\~na \textbf{Applications}
  \item Copiar la URL de JupyterHub (puerto 9443)
  \item Abrir en navegador, usuario: \texttt{jovyan} / pass: \texttt{jupyter}
  \item Subir el archivo \texttt{notebooks/agrosmart\_pipeline.ipynb}
  \item Men\'u: \textbf{Kernel} $\to$ \textbf{PySpark} $\to$ ejecutar todas las celdas
\end{enumerate}
\end{awsbox}

\subsection{Qu\'e hace el notebook (38 celdas)}

\begin{tabular}{lll}
\toprule
Secci\'on & Celdas & Acci\'on \\
\midrule
Setup SparkSession & 1--3 & Inicializar Spark con S3 como storage \\
Ingesta S3 Raw & 4--7 & Leer 3 CSV, revisar schema y nulos \\
Limpieza ETL & 8--13 & Dedup, imputaci\'on, tipado, fechas \\
Join e integraci\'on & 14--17 & Unir clima + IoT + producci\'on \\
\'Indices agroclim\'aticos & 18--22 & Calcular IHB, ITT, IEH, ET$_0$ \\
EDA & 23--26 & Estad\'isticas descriptivas, correlaci\'on \\
RandomForest (ML) & 27--30 & Clasificaci\'on riesgo tiz\'on, AUC=0.87 \\
GBT Regressor (ML) & 31--34 & Predicci\'on rendimiento, R$^2$=0.78 \\
Exportar Gold & 35--38 & Escribir 4 CSV Gold a S3 Gold \\
\bottomrule
\end{tabular}

\begin{okbox}[Archivos Gold generados en S3]
\begin{lstlisting}
s3://agrosmart-andino-demo/gold/lecturas_climaticas.csv     # 3,653 filas
s3://agrosmart-andino-demo/gold/alertas_activas.csv         # 500 filas
s3://agrosmart-andino-demo/gold/recomendaciones_riego.csv   # 300 filas
s3://agrosmart-andino-demo/gold/predicciones_rendimiento.csv # 6 filas
\end{lstlisting}
\end{okbox}

\begin{alertbox}[Descargar los CSV Gold para uso local]
Despu\'es de ejecutar el notebook, descargar los 4 CSV Gold a
\texttt{data/gold/} para poder usar el dashboard localmente sin EMR.
Desde AWS Console $\to$ S3 $\to$ gold/ $\to$ Download.
\end{alertbox}

\subsection{Ejecuci\'on local (alternativa sin EMR)}

\begin{localbox}[Si tienes PySpark instalado localmente]
\begin{lstlisting}[style=bash]
# Instalar PySpark
pip install pyspark

# Ejecutar el notebook localmente
cd agrosmart_andino
jupyter notebook notebooks/agrosmart_pipeline.ipynb
# NOTA: cambiar la ruta de S3 por la ruta local data/raw/
#   "s3://..." -> "data/raw/..."
\end{lstlisting}
\end{localbox}

% ============================================================
\section{Parte 5 -- Apache Cassandra (NoSQL)}
% ============================================================

\textbf{Objetivo:} Crear el esquema Cassandra y cargar los datos Gold.

\subsection{Opci\'on A: Cassandra local con Docker}

\begin{localbox}[Docker Desktop -- Instalar Cassandra 4.1]
\begin{lstlisting}[style=bash]
# Descargar imagen y ejecutar contenedor
docker pull cassandra:4.1
docker run --name cassandra-agrosmart \
  -p 9042:9042 \
  -e HEAP_NEWSIZE=128M \
  -e MAX_HEAP_SIZE=512M \
  -d cassandra:4.1

# Esperar ~30 segundos a que arranque
docker logs cassandra-agrosmart --tail 20
# Debe aparecer: "Starting listening for CQL clients on /0.0.0.0:9042"
\end{lstlisting}
\end{localbox}

\begin{localbox}[Crear el esquema CQL]
\begin{lstlisting}[style=bash]
# Copiar el schema al contenedor y ejecutar
docker cp scripts/cassandra/07_create_schema.cql cassandra-agrosmart:/tmp/

docker exec -it cassandra-agrosmart \
  cqlsh -f /tmp/07_create_schema.cql

# Verificar keyspace creado
docker exec -it cassandra-agrosmart \
  cqlsh -e "DESCRIBE KEYSPACES;"
# Debe aparecer: agrosmart_andino
\end{lstlisting}
\end{localbox}

\begin{localbox}[Cargar datos Gold en Cassandra]
\begin{lstlisting}[style=bash]
# Ejecutar desde agrosmart_andino/
# Modo dry-run (sin conexion real, solo valida datos):
python scripts/cassandra/08_load_data.py --dry-run

# Modo real (requiere Cassandra corriendo en localhost:9042):
python scripts/cassandra/08_load_data.py
\end{lstlisting}
\end{localbox}

\begin{okbox}[]
\begin{lstlisting}
[DRY-RUN] lecturas_climaticas  : 3,653 filas
[DRY-RUN] alertas_activas       : 500   filas
[DRY-RUN] recomendaciones_riego : 300   filas
[DRY-RUN] predicciones_rendimiento: 6   filas
Total: 4,459 filas procesadas en 0.1s (79,917 filas/seg)
\end{lstlisting}
\end{okbox}

\subsection{Opci\'on B: Cassandra en AWS Cloud9}

\begin{awsbox}[AWS Cloud9 -- Instalar Docker + Cassandra]
\begin{lstlisting}[style=bash]
# 1. Crear entorno Cloud9 en AWS Console
#    Tipo: EC2 t3.medium, Amazon Linux 2023

# 2. En la terminal de Cloud9, ejecutar el script:
bash scripts/cassandra/09_setup_cassandra_cloud9.sh

# El script automaticamente:
#   - Instala Docker
#   - Descarga cassandra:4.1
#   - Ejecuta el contenedor
#   - Crea el keyspace agrosmart_andino
#   - Crea las 5 tablas
\end{lstlisting}
\end{awsbox}

\begin{alertbox}[Compatibilidad Python 3.13 con cassandra-driver]
En Python 3.13 el m\'odulo \texttt{asyncore} fue eliminado. El script
\texttt{08\_load\_data.py} ya incluye el fix:
\begin{lstlisting}[style=python]
from cassandra.io.asyncioreactor import AsyncioConnection
cluster = Cluster(
    ['localhost'],
    connection_class=AsyncioConnection
)
\end{lstlisting}
No cambiar esto; es obligatorio para Python 3.13.
\end{alertbox}

\subsection{Estructura de tablas Cassandra}

\begin{tabular}{lll}
\toprule
Tabla & Partition Key & TTL \\
\midrule
\texttt{lecturas\_climaticas\_por\_dia} & (estacion, fecha DESC) & Sin TTL \\
\texttt{alertas\_activas} & (nodo\_id, fecha DESC) & 90 d\'ias \\
\texttt{recomendaciones\_riego} & (nodo\_id, fecha DESC) & 90 d\'ias \\
\texttt{predicciones\_rendimiento} & (anio, variedad) & Sin TTL \\
\texttt{estado\_nodos} & (nodo\_id) & Sin TTL \\
\bottomrule
\end{tabular}

% ============================================================
\section{Parte 6 -- Dashboard Streamlit}
% ============================================================

\textbf{Objetivo:} Ejecutar el dashboard interactivo de monitoreo en tiempo real.

\subsection{Prerrequisitos}

Antes de ejecutar el dashboard, verificar que existan los 4 CSV Gold en
\texttt{data/gold/}:

\begin{lstlisting}
data/gold/
+-- lecturas_climaticas.csv         (3,653 filas)
+-- alertas_activas.csv             (500 filas)
+-- recomendaciones_riego.csv       (300 filas)
+-- predicciones_rendimiento.csv    (6 filas)
\end{lstlisting}

Si no existen, copiarlos desde S3 o ejecutar las Partes 2--4 primero.

\subsection{Ejecutar el dashboard}

\begin{localbox}[IMPORTANTE: ejecutar DESDE la carpeta agrosmart\_andino/]
\begin{lstlisting}[style=bash]
# Navegar a la carpeta correcta
cd d:\Documents\GitHub\Maestria_IA_caso_practico\Big_data\trabajo_final\agrosmart_andino

# Ejecutar Streamlit
streamlit run dashboard/app.py --server.port 8501
\end{lstlisting}
\end{localbox}

\begin{alertbox}[Error frecuente: File does not exist: dashboard\textbackslash{}app.py]
\textbf{Causa:} El comando se ejecuta desde la carpeta equivocada.\\
\textbf{Soluci\'on:} Siempre ejecutar desde \texttt{agrosmart\_andino/}, no desde
\texttt{trabajo\_final/}.
\begin{lstlisting}[style=bash]
# INCORRECTO (desde trabajo_final/):
streamlit run dashboard/app.py

# CORRECTO (desde agrosmart_andino/):
cd agrosmart_andino
streamlit run dashboard/app.py --server.port 8501
\end{lstlisting}
\end{alertbox}

\begin{okbox}[Dashboard corriendo]
\begin{lstlisting}
  You can now view your Streamlit app in your browser.
  Local URL: http://localhost:8501
  Network URL: http://192.168.x.x:8501
\end{lstlisting}
Abrir \href{http://localhost:8501}{http://localhost:8501} en el navegador.
\end{okbox}

\subsection{Componentes del dashboard (11 KPIs)}

\begin{tabular}{llp{6cm}}
\toprule
\# & Componente & Descripci\'on \\
\midrule
1 & Sem\'aforo HTML & Verde/Naranja/Rojo seg\'un nivel de alerta \\
2 & KPI Nivel alerta & Texto coloreado con el estado del nodo \\
3 & KPI Humedad suelo & Porcentaje medio del nodo seleccionado \\
4 & KPI Recomendaci\'on & Texto de acci\'on de riego \\
5 & Tabla alertas & \'Ultimas 10 alertas por nodo \\
6 & M\'etricas 7 valores & T$_{\text{med}}$, T$_{\min}$, ET$_0$, d\'eficit, horas riego, etc. \\
7 & Temperatura 7 d\'ias & L\'inea m\'ax/med/m\'in con zona de helada \\
8 & Precipitaci\'on 30 d & Barras diarias \\
9 & Humedad suelo & L\'inea por nodo con bandas de estr\'es \\
10 & Heatmap ITT & Matriz sem$\times$d\'ia verde a rojo \\
11 & Gauge + barras GBT & Predicci\'on rendimiento 2025 por variedad \\
\bottomrule
\end{tabular}

\vspace{0.5em}
\textbf{Controles del sidebar:}
\begin{itemize}
  \item Selector de nodo: NODO-001 / NODO-002 / NODO-003
  \item Slider historial: 7 a 90 d\'ias
  \item Bot\'on \textit{Recargar datos}: fuerza recarga de CSV
  \item Refresco autom\'atico cada 5 minutos (\texttt{@st.cache\_data(ttl=300)})
\end{itemize}

% ============================================================
\section{Parte 7 -- Informe T\'ecnico en LaTeX}
% ============================================================

\textbf{Objetivo:} Compilar el informe t\'ecnico PDF de 20 p\'aginas.

\begin{localbox}[Compilar el informe (dos pasadas)]
\begin{lstlisting}[style=bash]
cd agrosmart_andino\docs

# Primera pasada (genera estructura)
pdflatex -interaction=nonstopmode informe_tecnico.tex

# Segunda pasada (resuelve referencias cruzadas e indice)
pdflatex -interaction=nonstopmode informe_tecnico.tex

# O usar el script batch incluido:
compilar.bat
\end{lstlisting}
\end{localbox}

\begin{alertbox}[Warning "Incompatible glue units" (x4)]
Este warning proviene de \texttt{tcolorbox} + \texttt{babel/spanish}. Es inofensivo;
el PDF se genera correctamente (20 p\'aginas). No es un error que impida la
compilaci\'on.
\end{alertbox}

\begin{okbox}[]
\texttt{docs/informe\_tecnico.pdf} --- 20 p\'aginas, $\approx$491 KB\\
Secciones: Portada, Resumen, 12 secciones t\'ecnicas, Referencias, Anexos A/B/C
\end{okbox}

% ============================================================
\section{Parte 8 -- Presentaci\'on Beamer}
% ============================================================

\textbf{Objetivo:} Compilar las diapositivas Beamer PDF de 17 slides.

\begin{localbox}[Compilar la presentaci\'on]
\begin{lstlisting}[style=bash]
cd agrosmart_andino\docs

# Dos pasadas para resolver el indice de diapositivas
pdflatex -interaction=nonstopmode presentacion.tex
pdflatex -interaction=nonstopmode presentacion.tex
\end{lstlisting}
\end{localbox}

\begin{okbox}[]
\texttt{docs/presentacion.pdf} --- 17 diapositivas, $\approx$495 KB\\
Duraci\'on estimada: 15 minutos + 5 minutos Q\&A
\end{okbox}

\subsection{Estructura de la presentaci\'on (guion 15 min)}

\begin{tabular}{clcp{5.5cm}}
\toprule
Slide & Secci\'on & Min & Puntos clave \\
\midrule
1 & Portada & 0:00 & T\'itulo, integrantes, docente \\
2 & \'Indice & 0:30 & Navegaci\'on de la presentaci\'on \\
3 & Problema & 1:30 & 3 riesgos, p\'erdida 38\%, oportunidad \\
4 & Objetivos & 3:00 & General + 5 espec\'ificos con m\'etricas \\
5 & Arquitectura & 4:30 & Flujo 5 capas, tabla tecnolog\'ias \\
6 & Datos & 6:00 & 72\,397 registros, 3 fuentes, validaci\'on \\
7 & ETL Spark & 7:30 & 3 fases, config EMR, tabla resultados \\
8 & \'Indices & 8:30 & F\'ormulas IHB/ITT/IEH/ET$_0$ \\
9 & ML & 10:00 & RF AUC=0.87, GBT R$^2$=0.78 \\
10 & Predicciones & 11:00 & 6 variedades, 9.2 t/ha 2025 \\
11 & Cassandra & 12:00 & 5 tablas, dry-run, design patterns \\
12 & Dashboard & 13:00 & 11 KPIs, ejecuci\'on, sem\'afor \\
13 & Sem\'aforos & 13:30 & Mock 3 nodos verde/naranja/rojo \\
14 & Resultados & 14:00 & Tabla global de m\'etricas \\
15 & Contribuciones & 14:30 & 4 aportes + limitaciones \\
16 & Conclusiones & 15:30 & 6 conclusiones enumeradas \\
17 & Cierre Q\&A & 16:00 & Gracias + preguntas \\
\bottomrule
\end{tabular}

% ============================================================
\section{Resumen: Comandos de Ejecuci\'on Completa}
% ============================================================

A continuaci\'on, el orden completo de ejecuci\'on para reproducir
todo el proyecto desde cero:

\begin{lstlisting}[style=bash]
# == PASO 0: Instalar dependencias ==========================
cd agrosmart_andino
pip install -r requirements.txt

# == PASO 1: Adquirir datos (Parte 2) =======================
python scripts/data_acquisition/01_download_nasa_power.py
python scripts/data_acquisition/02_generate_iot_data.py
python scripts/data_acquisition/03_prepare_midagri_data.py
python scripts/data_acquisition/04_validate_datasets.py

# == PASO 2: Subir a S3 (Parte 3) -- requiere AWS activo ====
# Copiar credenciales de AWS Academy en .env
python scripts/aws/05_upload_to_s3.py

# == PASO 3: ETL + ML en EMR (Parte 4) ======================
# Crear cluster EMR:
python scripts/aws/06_create_emr_cluster.py
# Abrir JupyterHub y ejecutar:
# notebooks/agrosmart_pipeline.ipynb
# Descargar 4 CSV Gold a data/gold/

# == PASO 4: Cassandra (Parte 5) ============================
docker run --name cassandra-agrosmart -p 9042:9042 \
  -e MAX_HEAP_SIZE=512M -d cassandra:4.1
# esperar 30 segundos...
docker cp scripts/cassandra/07_create_schema.cql cassandra-agrosmart:/tmp/
docker exec -it cassandra-agrosmart cqlsh -f /tmp/07_create_schema.cql
python scripts/cassandra/08_load_data.py --dry-run   # validar
python scripts/cassandra/08_load_data.py              # cargar real

# == PASO 5: Dashboard (Parte 6) ===========================
# IMPORTANTE: ejecutar desde agrosmart_andino/, no desde trabajo_final/
cd agrosmart_andino
streamlit run dashboard/app.py --server.port 8501
# Abrir: http://localhost:8501

# == PASO 6: Compilar documentacion (Partes 7 y 8) ==========
cd docs
pdflatex -interaction=nonstopmode informe_tecnico.tex
pdflatex -interaction=nonstopmode informe_tecnico.tex
pdflatex -interaction=nonstopmode presentacion.tex
pdflatex -interaction=nonstopmode presentacion.tex
\end{lstlisting}

% ============================================================
\section{Soluci\'on de Problemas Frecuentes}
% ============================================================

\begin{tabular}{p{4.5cm}p{5cm}p{4.5cm}}
\toprule
\textbf{Error} & \textbf{Causa} & \textbf{Soluci\'on} \\
\midrule
\texttt{File does not exist: dashboard\textbackslash{}app.py} &
Ejecutar desde carpeta incorrecta &
\texttt{cd agrosmart\_andino} antes de streamlit \\
\midrule
\texttt{NoCredentialsError} &
\texttt{.env} vac\'io o no encontrado &
Copiar credenciales de AWS Academy Details \\
\midrule
\texttt{asyncore} not found &
Python 3.13 elimin\'o asyncore &
Ya corregido: usar \texttt{AsyncioConnection} \\
\midrule
\texttt{applymap} deprecated &
Pandas $\ge$ 2.1 ren. \texttt{applymap} &
Ya corregido: usar \texttt{.style.map()} \\
\midrule
Incompatible glue units (LaTeX) &
tcolorbox + babel/spanish &
Es un warning inofensivo, el PDF se genera \\
\midrule
EMR cluster en estado \texttt{TERMINATED} &
Sesi\'on AWS de 3h expir\'o &
Reiniciar lab y crear nuevo cluster \\
\midrule
Cassandra: \texttt{Connection refused} &
Docker no iniciado &
\texttt{docker start cassandra-agrosmart} \\
\midrule
NASA POWER: timeout &
API lenta o red &
Ejecutar nuevamente (es gratuita, sin clave) \\
\bottomrule
\end{tabular}

% ============================================================
\section{Entregables Finales}
% ============================================================

\begin{tabular}{lll}
\toprule
Archivo & Descripci\'on & Tama\~no \\
\midrule
\texttt{docs/informe\_tecnico.pdf} & Informe t\'ecnico completo & 20 p\'ag., 491 KB \\
\texttt{docs/presentacion.pdf} & Diapositivas Beamer & 17 slides, 495 KB \\
\texttt{dashboard/app.py} & C\'odigo fuente dashboard & Streamlit \\
\texttt{notebooks/agrosmart\_pipeline.ipynb} & Notebook Spark EMR & 38 celdas \\
\texttt{data/gold/*.csv} & Datasets procesados & 4 archivos \\
\texttt{scripts/} & Todos los scripts (01--09) & 9 archivos \\
\bottomrule
\end{tabular}

\vspace{0.5em}
\begin{tcolorbox}[colback=agroclear!40, colframe=agrolight,
  title=M\'etricas clave para mencionar en la presentaci\'on,
  fonttitle=\bfseries\color{white}, colbacktitle=agrolight]
\begin{multicols}{2}
\begin{itemize}
  \item 72\,397 registros, 3 fuentes
  \item 3\,653 registros Gold integrados
  \item AUC-ROC = \textbf{0.87} (RandomForest)
  \item R$^2$ = \textbf{0.78} (GBT Regressor)
  \item 9.2 t/ha predicci\'on 2025
  \item 18 d\'ias helada / 47 tiz\'on / 31 h\'idrico
  \item 4\,459 filas Cassandra en 0.1 s
  \item $<$3 min ETL en EMR 6.15
  \item 11 componentes Streamlit
  \item $<$\$50/mes costo AWS estimado
\end{itemize}
\end{multicols}
\end{tcolorbox}

\end{document}
